\subsection{Indirect Support Prioritizes Drug Targets With Higher Clinical Success}

% STRUCTURE: This section has 4 paragraphs:
% P1: The gap (no non-genetic method predicts drug success) + continuous RS framework + weighted EA
% P2: Resources evaluated and curation
% P3: Main RS curve results (Figure 4a,c) + temporal holdout (Figure 4b)
% P4: Weighted EA by clinical area (Figure 5)
% Discussion-worthy deeper stories (oncology saturation, cardio ion channels, metabolic lipid/glucose)
% are flagged with DISCUSSION tags and written in the discussion section.

\begin{figure}[t!]
    \centering
    {{canva:5}}
    \caption{Drug target validation. (a) Continuous relative success (RS) as a function of the fraction of each resource prioritized ($\alpha$), comparing PIGEAN indirect, direct, and combined support (solid lines) to Open Targets Genetics (OTG) locus-to-gene share (dashed lines) and Open Targets Platform (OTP) combined and indirect support (dotted lines). Shaded bands denote 95\% confidence intervals computed by the delta method. Annotations indicate the number of target--indication (T-I) pairs prioritized at selected thresholds. (b) Temporal holdout: RS for drug targets launched after 2018, scored using PIGEAN with pre-2018 gene set libraries on 67 validation traits. (c) Cross-resource comparison of maximum RS achieved after at least 5\% of each resource is prioritized. For binary resources (Genebass, OMIM, Piccolo), RS is computed from a single supported/unsupported partition. Error bars denote 95\% delta method confidence intervals; numbers above each point indicate the number of T-I pairs in the supported set. Resources are grouped by curation source: Minikel curations (left), MAGMA (center), OTP (center-right), and PIGEAN (right), with score types distinguished by color (legend).}
    \label{fig:validation-pharma}
\end{figure}

% P1: The gap + continuous RS + weighted EA metrics
Minikel et al.\ demonstrated that drug targets with genetic support, defined as a gene--trait association in a genetic resource such as Open Targets Genetics (OTG), are approximately two-fold more likely to progress through clinical trials than unsupported targets \cite{minikel_refining_2024}. This finding has motivated efforts to use human genetic evidence for target selection, but genetic association data alone have limited coverage: even the largest compilation of OTG locus-to-gene (L2G) scores prioritizes only {{MINIKEL_OTG_ALL_N_TARGETS:,}} target--indication (T-I) pairs (Figure~\ref{fig:validation-pharma}c). Non-genetic evidence sources such as biological pathways, model organism phenotypes, and literature co-occurrence could substantially expand the number of evaluable targets. The Open Targets Platform (OTP) Europe PMC literature score is the closest existing non-genetic prioritization metric \cite{buniello_open_2025}, yet it achieves a maximum RS of only {{OTP_INDIRECT_RS:.2f}} (95\% CI: {{OTP_INDIRECT_RS_CI_LOWER:.2f}} to {{OTP_INDIRECT_RS_CI_UPPER:.2f}}) across {{OTP_INDIRECT_N_TARGETS:,}} T-I pairs (Figure~\ref{fig:validation-pharma}c), barely distinguishable from the null expectation of $\text{RS} = 1$. No method has demonstrated that non-genetic evidence can be systematically converted into a continuous score that meaningfully predicts drug target clinical success.

% TODO: CITE_OTP_EUROPEPMC - Add citation for Open Targets Platform Europe PMC score description

PIGEAN's indirect support fills this gap. Because PIGEAN, unlike binary classifiers, assigns continuous scores to gene--trait associations, we developed the \textit{continuous RS curve}, which traces $\text{RS}_{\text{cum}}$ as a function of the fraction $\alpha$ of T-I pairs prioritized (Section~\ref{sec:methods-rs}; Figure~\ref{fig:validation-pharma}a). This curve provides a complete profile of enrichment across all possible score thresholds, avoiding the need to select an arbitrary cutoff. To summarize the curve in a single scalar that balances prioritization quality against resource completeness, we defined the \textit{weighted excess area} (EA) above a baseline RS, computed as the integral of $(\text{RS}_{\text{cum}}(\alpha) - \text{RS}_{\text{baseline}})$ weighted by the number of T-I pairs at each threshold (Section~\ref{sec:methods-rs}). We report two variants: EA above $\text{RS} = 1$ ($\text{EA}|\text{RS}{=}1$), which captures total enrichment for successful targets across all thresholds, and EA above $\text{RS} = \text{Self}$ ($\text{EA}|\text{RS}{=}\text{Self}$), which uses the RS when the entire resource is designated as ``supported'' as the baseline, thereby isolating how much the resource's internal ranking adds beyond simple presence or absence in the resource. Because both metrics are weighted by target count, resources with greater coverage are rewarded: a resource that achieves high RS for many targets will score higher than one with equivalent RS for only a handful.

% P2: Resources evaluated and curation
We evaluated PIGEAN alongside several external resources using the continuous RS framework (Section~\ref{sec:methods-rs}). We matched gene--trait associations to T-I pairs from the Pharmaprojects clinical trial database across {{NUM_TOTAL_TRAITS_MOUSE_MSIGDB}} phenotypes using MeSH-based semantic similarity (threshold $\geq 0.8$). As benchmarks, we included (i) discrete curations by Minikel et al.\ \cite{minikel_refining_2024} from Genebass ($N$ = {{MINIKEL_GENEBASS_N_TARGETS}}), OMIM ($N$ = {{MINIKEL_OMIM_N_TARGETS}}), and Piccolo ($N$ = {{MINIKEL_PICCOLO_N_TARGETS}}), in which a T-I pair is either supported or not; (ii) continuous L2G share scores curated by Minikel et al.\ from OTG across five biobank resources (All OTG, FinnGen, GWAS Catalog, UKBB Neale, Saige); (iii) the Open Targets Platform (OTP) association scores, queried from the BigQuery public dataset filtering to non-drug evidence types (\texttt{datatypeId} $\in$ \{animal\_model, genetic\_association, literature\}), with the \textit{combined} score (harmonic mean of genetic and non-genetic evidence excluding drug-derived sources) and the \textit{indirect} score (Europe PMC literature evidence); (iv) MAGMA gene-level $p$-values \cite{de_leeuw_magma_2015} computed on the same summary statistics used for PIGEAN, restricted to genes with $p < 2.5 \times 10^{-6}$; and (v) PIGEAN posterior probabilities (combined, direct, and indirect) from the Mouse + MSigDB model. For continuous resources in Figure~\ref{fig:validation-pharma}c, the reported RS is the maximum achieved after at least 5\% of the resource is prioritized, ensuring numerical stability; for discrete resources, RS is computed from the single supported/unsupported partition.

% P3: Main RS results
% P3a: Indirect support works -- using genetics as anchor to leverage non-genetic connections
Across all sources, PIGEAN indirect support achieved RS = {{PIGEAN_INDIRECT_RS_ALL:.2f}} (95\% CI: {{PIGEAN_INDIRECT_RS_ALL_CI_LOWER:.2f}} to {{PIGEAN_INDIRECT_RS_ALL_CI_UPPER:.2f}}) while prioritizing {{PIGEAN_INDIRECT_N_TARGETS_ALL:,}} T-I pairs, covering 78\% of OTP indirect's target set ({{OTP_INDIRECT_N_TARGETS:,}}) at 1.8-fold higher RS ({{PIGEAN_INDIRECT_RS_ALL:.2f}} versus {{OTP_INDIRECT_RS:.2f}}; Figure~\ref{fig:validation-pharma}c). PIGEAN combined support (direct plus indirect) achieved RS = {{PIGEAN_COMBINED_RS_ALL:.2f}} (95\% CI: {{PIGEAN_COMBINED_RS_ALL_CI_LOWER:.2f}} to {{PIGEAN_COMBINED_RS_ALL_CI_UPPER:.2f}}; $N$ = {{PIGEAN_COMBINED_N_TARGETS_ALL:,}}), and PIGEAN direct support achieved RS = {{PIGEAN_DIRECT_RS_ALL:.2f}} (95\% CI: {{PIGEAN_DIRECT_RS_ALL_CI_LOWER:.2f}} to {{PIGEAN_DIRECT_RS_ALL_CI_UPPER:.2f}}; $N$ = {{PIGEAN_DIRECT_N_TARGETS_ALL:,}}). Indirect support (RS = {{PIGEAN_INDIRECT_RS_ALL:.2f}}) approached combined support (RS = {{PIGEAN_COMBINED_RS_ALL:.2f}}), indicating that by using genetic associations as an anchor, PIGEAN can leverage shared pathway and annotation structure to prioritize a substantially larger set of targets than genetics alone, at comparable enrichment levels. The indirect score requires direct genetic evidence to calibrate which annotations are trait-relevant (Section~\ref{sec:methods-rs}), but the resulting annotation-based prioritization extends far beyond the directly associated genes, reaching {{PIGEAN_INDIRECT_N_TARGETS_ALL:,}} T-I pairs compared with {{MINIKEL_OTG_ALL_N_TARGETS:,}} from OTG L2G.

% P3b: Comparison to curated resources -- more targets at comparable RS
The continuous RS curve (Figure~\ref{fig:validation-pharma}a) reveals a fundamental tradeoff between enrichment magnitude and target coverage. Among the Minikel curations, the highest RS values were achieved by the smallest, most stringent resources: UKBB Neale L2G share (RS = {{MINIKEL_UKBB_RS:.2f}}; 95\% CI: {{MINIKEL_UKBB_RS_CI_LOWER:.2f}} to {{MINIKEL_UKBB_RS_CI_UPPER:.2f}}; $N$ = {{MINIKEL_UKBB_N_TARGETS}}), OMIM (RS = {{MINIKEL_OMIM_RS:.2f}}; 95\% CI: {{MINIKEL_OMIM_RS_CI_LOWER:.2f}} to {{MINIKEL_OMIM_RS_CI_UPPER:.2f}}; $N$ = {{MINIKEL_OMIM_N_TARGETS}}), and Saige (RS = {{MINIKEL_SAIGE_RS:.2f}}; 95\% CI: {{MINIKEL_SAIGE_RS_CI_LOWER:.2f}} to {{MINIKEL_SAIGE_RS_CI_UPPER:.2f}}; $N$ = {{MINIKEL_SAIGE_N_TARGETS}}). MAGMA achieved RS = {{MAGMA_RS:.2f}} (95\% CI: {{MAGMA_RS_CI_LOWER:.2f}} to {{MAGMA_RS_CI_UPPER:.2f}}) for {{MAGMA_N_TARGETS}} T-I pairs. These resources achieve high RS but for target sets that are one to two orders of magnitude smaller than PIGEAN's: PIGEAN combined support covers 63-fold more T-I pairs than UKBB Neale ({{PIGEAN_COMBINED_N_TARGETS_ALL:,}} versus {{MINIKEL_UKBB_N_TARGETS}}) and 12-fold more than All OTG L2G share ({{PIGEAN_COMBINED_N_TARGETS_ALL:,}} versus {{MINIKEL_OTG_ALL_N_TARGETS:,}}), at only modestly lower RS ({{PIGEAN_COMBINED_RS_ALL:.2f}} versus {{MINIKEL_UKBB_RS:.2f}} and {{MINIKEL_OTG_ALL_RS:.2f}}, respectively). Notably, the OTG L2G curve in Figure~\ref{fig:validation-pharma}a shows relatively flat RS across L2G scores, such that even scores approaching 1.0 do not yield substantially higher RS than scores near 0.1, consistent with the findings of Minikel et al.\ \cite{minikel_refining_2024}. By contrast, PIGEAN's RS curve rises steeply at stringent thresholds (Figure~\ref{fig:validation-pharma}a), indicating that its scoring has greater discriminative power for identifying clinically successful targets.

% P3c: OTP fails to rank targets
OTP's combined and indirect scores, despite covering {{OTP_COMBINED_N_TARGETS:,}} and {{OTP_INDIRECT_N_TARGETS:,}} T-I pairs respectively, achieved RS of only {{OTP_COMBINED_RS:.2f}} (95\% CI: {{OTP_COMBINED_RS_CI_LOWER:.2f}} to {{OTP_COMBINED_RS_CI_UPPER:.2f}}) and {{OTP_INDIRECT_RS:.2f}} (95\% CI: {{OTP_INDIRECT_RS_CI_LOWER:.2f}} to {{OTP_INDIRECT_RS_CI_UPPER:.2f}}), hovering near the null across the full range of thresholds (Figure~\ref{fig:validation-pharma}a, dotted lines). The baseline RS (i.e., when the entire resource is designated as supported) was {{OTP_COMBINED_BASELINE_RS:.2f}} for OTP combined and {{OTP_INDIRECT_BASELINE_RS:.2f}} for OTP indirect, both at or below 1.0, indicating that OTP's full target set is not enriched for clinical success relative to the Pharmaprojects background. This stands in contrast to PIGEAN, where the baseline RS of {{PIGEAN_COMBINED_BASELINE_RS_ALL:.2f}} indicates mild enrichment even before applying any score threshold.

% P3d: Source-matched comparisons
Source-matched comparisons confirm that PIGEAN's advantage is not driven by trait selection. Restricting to GWAS Catalog (GCAT) traits, PIGEAN direct support achieved RS = {{PIGEAN_DIRECT_RS_GCAT:.2f}} ($N$ = {{PIGEAN_DIRECT_N_TARGETS_GCAT:,}}), compared with {{PIGEAN_COMBINED_RS_GCAT:.2f}} for combined ($N$ = {{PIGEAN_COMBINED_N_TARGETS_GCAT:,}}) and {{PIGEAN_INDIRECT_RS_GCAT:.2f}} for indirect ($N$ = {{PIGEAN_INDIRECT_N_TARGETS_GCAT:,}}). For portal traits, indirect and combined support were nearly identical (RS = {{PIGEAN_INDIRECT_RS_PORTAL:.2f}} and {{PIGEAN_COMBINED_RS_PORTAL:.2f}}, respectively; both $N$ = {{PIGEAN_COMBINED_N_TARGETS_PORTAL:,}}), suggesting that for common diseases with full summary statistics, indirect evidence captures nearly all of the signal present in combined support. For Orphanet rare diseases, indirect support achieved the highest RS of any source--trait combination: RS = {{PIGEAN_INDIRECT_RS_ORPHANET:.2f}} (95\% CI: {{PIGEAN_INDIRECT_RS_ORPHANET_CI_LOWER:.2f}} to {{PIGEAN_INDIRECT_RS_ORPHANET_CI_UPPER:.2f}}; $N$ = {{PIGEAN_INDIRECT_N_TARGETS_ORPHANET:,}}), essentially equivalent to combined support (RS = {{PIGEAN_COMBINED_RS_ORPHANET:.2f}}; $N$ = {{PIGEAN_COMBINED_N_TARGETS_ORPHANET:,}}). This result indicates that indirect support is especially effective for rare diseases, where pathway biology is well-characterized but direct genetic evidence may be limited to small numbers of known causal genes.

% P3e: Temporal holdout
To assess whether the observed enrichment is driven by circularity, specifically gene set annotations that encode existing drug target knowledge, we performed a temporal holdout analysis (Section~\ref{sec:methods-rs}; Figure~\ref{fig:validation-pharma}b). PIGEAN was re-run using only gene set libraries curated from pre-2018 data, and evaluation was restricted to T-I pairs for drugs launched after 2018 across 67 validation traits. Under this temporal holdout, indirect support achieved a maximum RS of {{MAX_RS_INDIRECT_TEMPORAL_HOLDOUT:.2f}} (95\% CI: {{MAX_RS_INDIRECT_TEMPORAL_HOLDOUT_CI_LOWER:.2f}} to {{MAX_RS_INDIRECT_TEMPORAL_HOLDOUT_CI_UPPER:.2f}}; $N$ = {{NUM_DRUG_TARGETS_SUPPORTED_INDIRECT_TEMPORAL_HOLDOUT}}) and combined support achieved {{MAX_RS_COMBINED_TEMPORAL_HOLDOUT:.2f}} (95\% CI: {{MAX_RS_COMBINED_TEMPORAL_HOLDOUT_CI_LOWER:.2f}} to {{MAX_RS_COMBINED_TEMPORAL_HOLDOUT_CI_UPPER:.2f}}; $N$ = {{NUM_DRUG_TARGETS_SUPPORTED_COMBINED_TEMPORAL_HOLDOUT}}). At 2\% of the resource, RS stabilized at {{RS_INDIRECT_TEMPORAL_HOLDOUT_STABLE_AT_ALPHA_002:.2f}} for indirect support ($N$ = {{NUM_TARGETS_INDIRECT_TEMPORAL_HOLDOUT_AT_ALPHA_002}}) and {{RS_COMBINED_TEMPORAL_HOLDOUT_STABLE_AT_ALPHA_002:.2f}} for combined support ($N$ = {{NUM_TARGETS_COMBINED_TEMPORAL_HOLDOUT_AT_ALPHA_002}}), confirming that the enrichment is not an artifact of circularity with existing drug target annotations.

\begin{figure}[t!]
    \centering
    {{canva:6}}
    \caption{Weighted excess area (EA) above RS baseline across therapeutic areas. Each cell shows the weighted EA (top) and the number of supported T-I pairs (bottom, in parentheses) for a given method and therapeutic area. Left panel: baseline $\text{RS} = 1$ (comparison against all targets). Right panel: baseline $\text{RS} = \text{Self}$ (comparison against each method's own unsupported targets). Bold values indicate the highest weighted EA within each therapeutic area. Color intensity is scaled per column (therapeutic area) to highlight the best-performing method within each area.}
    \label{fig:weighted-ea}
\end{figure}

% P4: Weighted EA by clinical area
% P4a: Overall weighted EA
To compare resources in a manner that jointly accounts for enrichment magnitude and target coverage, we computed the weighted EA across 17 therapeutic areas (Figure~\ref{fig:weighted-ea}). Across all areas combined, PIGEAN combined support achieved the highest $\text{EA}|\text{RS}{=}1$ ({{WEA_PIGEAN_COMBINED_ALL:,}}; $N$ = {{PIGEAN_COMBINED_N_TARGETS_ALL:,}}), followed by PIGEAN indirect support ({{WEA_PIGEAN_INDIRECT_ALL:,}}; $N$ = {{PIGEAN_INDIRECT_N_TARGETS_ALL:,}}), PIGEAN direct support ({{WEA_PIGEAN_DIRECT_ALL:,}}; $N$ = {{PIGEAN_DIRECT_N_TARGETS_ALL:,}}), OTP combined ({{WEA_OTP_COMBINED_ALL:,}}; $N$ = {{OTP_COMBINED_N_TARGETS:,}}), OTP indirect ({{WEA_OTP_INDIRECT_ALL:,}}; $N$ = {{OTP_INDIRECT_N_TARGETS:,}}), OTG L2G share ({{WEA_OTG_L2G_ALL:,}}; $N$ = {{MINIKEL_OTG_ALL_N_TARGETS:,}}), and MAGMA ({{WEA_MAGMA_ALL}}; $N$ = {{MAGMA_N_TARGETS}}). Notably, PIGEAN combined and indirect support dominated despite covering fewer T-I pairs than OTP ({{PIGEAN_COVERAGE_FRACTION:.1%}} versus {{OTP_COVERAGE_FRACTION:.1%}} of all T-I pairs): because the weighted EA metric rewards coverage, OTP's broader target set should confer an advantage, yet PIGEAN's per-target enrichment signal was strong enough to more than compensate. On the $\text{EA}|\text{RS}{=}\text{Self}$ metric, which isolates ranking ability from baseline target quality, PIGEAN combined support again led ({{WEA_SELF_PIGEAN_COMBINED_ALL:,}}), followed by PIGEAN indirect ({{WEA_SELF_PIGEAN_INDIRECT_ALL:,}}), OTP combined ({{WEA_SELF_OTP_COMBINED_ALL:,}}), OTG L2G share ({{WEA_SELF_OTG_L2G_ALL}}), and MAGMA ({{WEA_SELF_MAGMA_ALL}}). This confirms that PIGEAN's scoring provides genuine ranking information, not merely a better pool of targets, and that indirect support alone captures the majority of this discriminative signal.

% P4b: Areas where PIGEAN dominates
The per-area breakdown reveals that PIGEAN's advantage is concentrated in therapeutic areas where pathway biology provides strong signal (Figure~\ref{fig:weighted-ea}). In oncology, PIGEAN indirect support achieved $\text{EA}|\text{RS}{=}1$ = {{WEA_PIGEAN_INDIRECT_ONCO:,}} ($N$ = {{WEA_PIGEAN_INDIRECT_ONCO_N:,}}), while OTP combined was negative ($\text{EA}|\text{RS}{=}1$ = {{WEA_OTP_COMBINED_ONCO:,}}; $N$ = {{WEA_OTP_COMBINED_ONCO_N:,}}), indicating that OTP's scoring actively anti-predicts oncology drug success. In neurology, PIGEAN combined achieved $\text{EA}|\text{RS}{=}1$ = {{WEA_PIGEAN_COMBINED_NEURO:,}} ($N$ = {{WEA_PIGEAN_COMBINED_NEURO_N:,}}) and PIGEAN indirect achieved {{WEA_PIGEAN_INDIRECT_NEURO:,}}, compared with {{WEA_OTP_COMBINED_NEURO}} for OTP combined ($N$ = {{WEA_OTP_COMBINED_NEURO_N:,}}), a six-fold advantage. In cardiovascular disease, PIGEAN indirect support ($\text{EA}|\text{RS}{=}1$ = {{WEA_PIGEAN_INDIRECT_CARDIO}}; $N$ = {{WEA_PIGEAN_INDIRECT_CARDIO_N:,}}) outperformed PIGEAN combined support ({{WEA_PIGEAN_COMBINED_CARDIO}}; $N$ = {{WEA_PIGEAN_COMBINED_CARDIO_N:,}}), one of the few areas where adding direct genetic evidence reduced performance. This suggests that pathway biology, including ion channels, contractile machinery, and lipid metabolism, is a stronger predictor of cardiovascular drug success than direct GWAS association. In endocrine disorders, PIGEAN indirect support likewise led ($\text{EA}|\text{RS}{=}1$ = {{WEA_PIGEAN_INDIRECT_ENDO}}; $N$ = {{WEA_PIGEAN_INDIRECT_ENDO_N}}).

% P4c: Areas where OTP is competitive
OTP performed competitively in therapeutic areas with well-characterized druggable landscapes and deep literature coverage. In immunology, OTP indirect support achieved $\text{EA}|\text{RS}{=}1$ = {{WEA_OTP_INDIRECT_IMMUNE}} ($N$ = {{WEA_OTP_INDIRECT_IMMUNE_N:,}}), far exceeding PIGEAN indirect support ({{WEA_PIGEAN_INDIRECT_IMMUNE}}; $N$ = {{WEA_PIGEAN_INDIRECT_IMMUNE_N:,}}). In metabolic disease, OTP combined ($\text{EA}|\text{RS}{=}1$ = {{WEA_OTP_COMBINED_METABOLIC:,}}; $N$ = {{WEA_OTP_COMBINED_METABOLIC_N}}) was narrowly ahead of PIGEAN combined ({{WEA_PIGEAN_COMBINED_METABOLIC:,}}; $N$ = {{WEA_PIGEAN_COMBINED_METABOLIC_N:,}}), though the $\text{EA}|\text{RS}{=}\text{Self}$ metric reversed this: PIGEAN combined ({{WEA_SELF_PIGEAN_COMBINED_METABOLIC:,}}) substantially outperformed OTP combined ({{WEA_SELF_OTP_COMBINED_METABOLIC}}), indicating that OTP has a strong pool of metabolic targets but PIGEAN ranks them more effectively. OTP also led in dermatology ($\text{EA}|\text{RS}{=}1$ = {{WEA_OTP_INDIRECT_DERM}}; $N$ = {{WEA_OTP_INDIRECT_DERM_N:,}}) and musculoskeletal disease ({{WEA_OTP_COMBINED_MUSCULO}}; $N$ = {{WEA_OTP_COMBINED_MUSCULO_N}}). In infection, all methods achieved negative or near-zero weighted EA (Figure~\ref{fig:weighted-ea}), consistent with the expectation that gene prioritization approaches are less applicable to infectious disease, where drug targets are often pathogen-derived rather than human genes.

% DISCUSSION: The oncology, neurology, cardiovascular, immune, and metabolic stories
% contain deeper mechanistic analysis that belongs in the Discussion section.
